\begin{algorithm}[htbp]
\caption{Skip Module Hyperparameter Selection}
\label{alg:hyperparam}
\begin{algorithmic}[1]

\Require Parameter space $\Theta$, validation set $D_{\text{val}}$,
DL model $\mathcal{M}$, recall tolerance $\delta_R$,
overhead ratio $\beta$, weights $w_S, w_R$
\Ensure Optimal parameters $\theta^*$

\State baseline $\gets$ \textsc{EvaluatePipeline}($D_{\text{val}}$, skip=None, $\mathcal{M}$)
\State Extract $R_{\text{base}}$, $T_{\text{DL}}$ from baseline
\State best\_score $\gets -\infty$; \quad $\theta^* \gets$ None

\For{each $\theta \in \Theta$}
\State results $\gets$ \textsc{EvaluatePipeline}($D_{\text{val}}$, $\mathcal{S}(\theta)$, $\mathcal{M}$)
\State Extract $R(\theta)$, $S_r(\theta)$, $T_{\text{skip}}(\theta)$ from results
\If{$R(\theta) \geq R_{\text{base}} - \delta_R$
\textbf{ and }
$T_{\text{skip}}(\theta) \leq \beta \cdot T_{\text{DL}}$}
\State $\rho(\theta) \gets 1 - \dfrac{R_{\text{base}} - R(\theta)}{\delta_R}$
\State score $\gets w_S\, S_r(\theta) + w_R\, \rho(\theta)$
\If{score $>$ best\_score}
\State best\_score $\gets$ score
\State $\theta^* \gets \theta$
\EndIf
\EndIf
\EndFor

\State \Return $\theta^*$

\end{algorithmic}
\end{algorithm}

The skip module $\mathcal{S}$ contains several hyperparameters (e.g., motion
threshold, accumulation window size) that can significantly impact the
performance of the overall system. To identify the optimal configuration, we
employ a grid search based hyperparameter optimization procedure on a validation
set $D_{\text{val}}$ derived from our video dataset with respect to several
system-wide metrics defined in Section \ref{sec:metrics}. More specifically, to
select the optimal hyperparameters for the proposed skip module, we employ a
constrained optimization procedure on the validation set $D_{\text{val}}$. Let $R_{\text{base}}$ denote the baseline recall, i.e., the fraction of fire/smoke frames correctly detected by the pipeline without skipping, and let $T_{\text{DL}}$ denote its mean per-frame inference time. For
each candidate parameter set $\theta \in \Theta$ obtained by grid search, we
evaluate the full pipeline $\text{Read} \rightarrow \mathcal{S}(\theta)
\rightarrow [\mathcal{M}]$ and compute three metrics: recall $R(\theta)$ (the fraction of fire/smoke frames correctly detected), skip rate $S_r(\theta)$ (the fraction of ground-truth negative frames skipped by $\mathcal{S}$), and mean per-frame skip module time $T_{\text{skip}}(\theta)$. Formal metric definitions are given in Section~\ref{sec:metrics}.


The skip rate $S_r(\theta)$ is computed on $D_{\text{val}}$ following the
definition in Section \ref{sec:metrics}, with $N_{\text{neg}}$ instantiated as
$N_{\text{neg}}^{\text{val}} = |\{f_i \in D_{\text{val}} : y_i = 0\}|$, which is
fixed by the validation set and independent of $\theta$.

The mean per-frame inference times are defined as

\[ T_{\text{DL}} = \frac{1}{|D_{\text{val}}|} \sum_{i=1}^{|D_{\text{val}}|}
t_{\text{DL}}(f_i), \qquad T_{\text{skip}}(\theta) = \frac{1}{|D_{\text{val}}|}
\sum_{i=1}^{|D_{\text{val}}|} t_{\text{skip}}(f_i, \theta), \]

where $f_i$ denotes the $i$-th frame in $D_{\text{val}}$, and both averages are
taken over all frames to reflect the runtime cost on a typical video stream.

\textbf{Feasibility constraints.} A candidate $\theta$ is considered feasible
only if it satisfies two hard constraints:

\[ R(\theta) \geq R_{\text{base}} - \delta_R, \qquad T_{\text{skip}}(\theta)
\leq \beta \cdot T_{\text{DL}}, \]

where recall tolerance $\delta_R > 0$ is the maximum allowable absolute recall
drop, and overhead ratio $\beta \in (0, 1)$ bounds the skip module overhead as a fraction of one full DL
inference. For example, setting $\delta_R = 0.01$ and $\beta = 0.10$ means the
system tolerates at most a 1\% recall drop and requires the skip module to run
within 10\% of the cost of a single DL inference.

\textbf{Scoring feasible candidates.} Among feasible candidates, the only
term that requires explicit optimization is recall retention, defined as

\[ \rho(\theta) = 1 - \frac{R_{\text{base}} - R(\theta)}{\delta_R}, \]

which equals $1$ when recall matches the baseline and $0$ when recall reaches
the lowest acceptable level $R_{\text{base}} - \delta_R$. This term is bounded
in $[0, 1]$ over the feasible region, making it directly comparable to
$S_r(\theta)$ in the weighted score.

\textbf{Optimal selection.} The optimal parameter set $\theta^*$ is selected as

\[ \theta^* = \arg\max_{\theta \in \Theta_{\text{feasible}}} \left[ w_R  \rho(\theta) + w_S S_r(\theta) \right], \]

where $\Theta_{\text{feasible}} = \{\theta \in \Theta : R(\theta) \geq
R_{\text{base}} - \delta_R \text{ and } T_{\text{skip}}(\theta) \leq \beta \cdot
T_{\text{DL}}\}$, and the nonnegative weights satisfy $w_S + w_R = 1$.

\textbf{Theoretical justification}: The skip module $\mathcal{S}$ acts as a
conservative gate: skipped frames are labeled negative directly, while passed
frames are processed by the same downstream detector $\mathcal{M}$ as in the
baseline. Therefore, every false alarm produced by the skip-enabled system is
also present in the baseline, implying $\mathrm{FAR}(\theta) \leq \mathrm{FAR}_{\text{base}}$, where $\mathrm{FAR}(\theta)$ and $\mathrm{FAR}_{\text{base}}$ denote the false alarm rates of the skip-enabled and baseline systems, respectively.The primary safety risk is recall degradation
caused by incorrectly skipped fire/smoke frames, which motivates enforcing the
recall constraint as a hard gate prior to any candidate ranking.

% ! Start justification here

\textbf{Theoretical justification for excluding FAR from the scoring objective.}
One might consider including a normalized FAR-reduction term

\[ F(\theta) = \frac{\mathrm{FAR}_{\text{base}} -
\mathrm{FAR}(\theta)}{\mathrm{FAR}_{\text{base}}}, \]

which measures the relative reduction in false alarm rate with respect to the
baseline, in the scoring objective. We argue that doing so is both structurally
unsound and statistically redundant, and therefore exclude it.

\textit{Structural argument.}
The skip module $\mathcal{S}$ labels every skipped frame as negative
($\hat{y}_t = 0$), so a skipped frame can never generate a new false positive.
Consequently, $\mathrm{FAR}(\theta) \leq \mathrm{FAR}_{\text{base}}$ holds for
all $\theta$, implying $F(\theta) \geq 0$ unconditionally. A scoring term that
is non-negative for every feasible candidate provides no discriminating signal:
it uniformly rewards all candidates and cannot distinguish configurations that
skip more effectively from those that do not. FAR reduction is therefore a
guaranteed structural side-effect of the skip mechanism, not an independently
optimizable objective.

\textit{Statistical argument.}
Let $N_{\text{neg}}$ be the total number of ground-truth negative frames in
$D_{\text{val}}$, $FP_{\text{base}}$ the number of those frames the baseline
classifier $\mathcal{M}$ labels as false positives, and
$\mathrm{FAR}_{\text{base}} = FP_{\text{base}} / N_{\text{neg}}$.
Among the $N_{\text{skip\_neg}}$ negative frames skipped by $\mathcal{S}(\theta)$,
let $a$ be the count that $\mathcal{M}$ \emph{would have} labeled as false
positives and $b$ the count it would have labeled correctly, so that
$a + b = N_{\text{skip\_neg}}$.
Since skipped frames are labeled $0$ and cannot be false positives,
$FP(\theta) = FP_{\text{base}} - a$, and substituting into $F(\theta)$ gives:

\[ F(\theta)
= \frac{a / N_{\text{neg}}}{FP_{\text{base}} / N_{\text{neg}}}
= \frac{a}{FP_{\text{base}}},
\qquad \text{while} \qquad
S_r(\theta) = \frac{a + b}{N_{\text{neg}}} \]

Both terms share variable $a$, so the question is whether $\mathcal{S}$ tends
to skip group-$a$ frames more often than group-$b$ frames.
Since $\mathcal{S}$ operates solely on inter-frame pixel motion and has no
access to $\mathcal{M}$'s weights or outputs, its skip decision is
\emph{statistically independent} of the FP/TN label.
Under this independence, the skip module samples blindly from the negative pool,
so each sub-type appears in proportion to its share in the full pool.
The expected count of group-$a$ frames among the skipped set is therefore:

\[ \mathbb{E}[a]
= N_{\text{skip\_neg}} \times \mathrm{FAR}_{\text{base}}
= N_{\text{skip\_neg}} \times \frac{FP_{\text{base}}}{N_{\text{neg}}} \]

Since $FP_{\text{base}}$ is a fixed constant, it passes through the expectation
unchanged, and $F(\theta)$ evaluates to:

\[ \mathbb{E}[F(\theta)]
= \frac{\mathbb{E}[a]}{FP_{\text{base}}}
= \frac{N_{\text{skip\_neg}} \times \dfrac{FP_{\text{base}}}{N_{\text{neg}}}}{FP_{\text{base}}}
= \frac{N_{\text{skip\_neg}}}{N_{\text{neg}}}
= S_r(\theta) \]

The $FP_{\text{base}}$ terms cancel exactly. In expectation, $F(\theta)$
is therefore identical to $S_r(\theta)$ and carries no independent information.
Including both in the scoring objective would silently assign skip-rate
behaviour a combined weight of $w_S + w_F$. We therefore exclude $F(\theta)$
from the scoring criterion and retain $\mathrm{FAR}(\theta)$ solely as an
observed consequence metric in the results tables.


% ! A long justification version
% \textbf{Theoretical justification for excluding FAR from the scoring objective.}
% One might consider including a normalized FAR-reduction term

% \[ F(\theta) = \frac{\mathrm{FAR}_{\text{base}} -
% \mathrm{FAR}(\theta)}{\mathrm{FAR}_{\text{base}}} \]

% which measures the relative reduction in false alarm rate with respect to the
% baseline, in the scoring objective. We argue that doing so is both structurally
% unsound and statistically redundant, and therefore exclude it.

% \textit{Structural argument.}
% The skip module $\mathcal{S}$ labels every skipped frame as negative
% ($\hat{y}_t = 0$). Since a false positive requires the system to output $1$ on a
% ground-truth negative frame, a skipped frame can never generate a new false
% positive. Consequently, $\mathrm{FAR}(\theta) \leq \mathrm{FAR}_{\text{base}}$
% holds for all $\theta$, which implies $F(\theta) \geq 0$ unconditionally.
% A scoring term that is non-negative for every feasible candidate provides no
% discriminating signal: it uniformly rewards all candidates and cannot distinguish
% configurations that skip more effectively from those that do not. FAR reduction
% is therefore a guaranteed structural side-effect of the skip mechanism, not an
% independently optimizable objective.

% \textit{Statistical argument.}
% We show that $F(\theta)$ is statistically redundant with respect to $S_r(\theta)$,
% meaning they carry the same information in expectation, and including both in the
% scoring objective is therefore unnecessary.

% Let $N_{\text{neg}}$ denote the total number of ground-truth negative
% (background-only) frames in the validation set $D_{\text{val}}$, and let
% $FP_{\text{base}}$ denote the number of those frames that the baseline classifier
% $\mathcal{M}$ incorrectly labels as fire or smoke --- its false positives.
% The baseline false alarm rate is then
% $\mathrm{FAR}_{\text{base}} = FP_{\text{base}} / N_{\text{neg}}$,
% which is the fraction of all negative frames that the classifier gets wrong.

% Now consider a candidate skip configuration $\theta$.
% Let $N_{\text{skip\_neg}}$ denote the number of negative frames that
% $\mathcal{S}(\theta)$ skips --- that is, frames assigned $s_t = 0$ without
% invoking $\mathcal{M}$.
% Among these $N_{\text{skip\_neg}}$ skipped frames, let $a$ denote the count of
% frames that $\mathcal{M}$ \emph{would have} labeled as false positives had they
% not been skipped, and let $b$ denote the count of frames it would have labeled
% correctly as true negatives, so that $a + b = N_{\text{skip\_neg}}$.

% Since every skipped frame is assigned label $0$ and therefore cannot be a false
% positive, the false positive count of the skip-enabled system satisfies
% $FP(\theta) = FP_{\text{base}} - a$, which gives
% $\mathrm{FAR}_{\text{base}} - \mathrm{FAR}(\theta)
% = (FP_{\text{base}} - FP(\theta)) / N_{\text{neg}}
% = a / N_{\text{neg}}$.
% Substituting into the definition of $F(\theta)$ yields:

% \[ F(\theta)
% = \frac{\mathrm{FAR}_{\text{base}} - \mathrm{FAR}(\theta)}{\mathrm{FAR}_{\text{base}}}
% = \frac{a / N_{\text{neg}}}{FP_{\text{base}} / N_{\text{neg}}}
% = \frac{a}{FP_{\text{base}}} \]

% Using the same counts, the skip rate $S_r(\theta)$ --- the fraction of all
% negative frames that were skipped --- expands as:

% \[ S_r(\theta) = \frac{N_{\text{skip\_neg}}}{N_{\text{neg}}} = \frac{a + b}{N_{\text{neg}}} \]

% The two terms therefore share variable $a$ but differ in what they divide by:
% $S_r$ counts all skipped negative frames $(a + b)$ over the full negative pool
% $N_{\text{neg}}$, while $F$ counts only the prevented false positives $(a)$ over
% the baseline FP count $FP_{\text{base}}$.
% The central question is whether the skip module has any systematic tendency to
% target frames in group $a$ over group $b$.

% The skip module $\mathcal{S}$ makes its decision based solely on inter-frame
% pixel motion: it compares the current frame against the previous one and skips
% if no significant motion is detected.
% It has no access to $\mathcal{M}$'s internal weights, learned feature
% representations, or prediction outputs.
% The classifier's false positives arise from background scenes that
% \emph{visually resemble} fire or smoke to $\mathcal{M}$ --- a property encoded
% entirely within the classifier's learned parameters, which is invisible to a
% motion detector.
% Therefore, the skip decision is \emph{statistically independent} of the FP/TN
% label: knowing that a frame was skipped provides no information about whether
% $\mathcal{M}$ would have misclassified it, and vice versa.

% Because of this independence, the skip module draws the $N_{\text{skip\_neg}}$
% skipped frames from the negative pool without any preference for one label over
% the other --- it samples blindly with respect to the FP/TN property.
% When items are sampled blindly from a pool, each sub-type appears in the sample
% in proportion to its share in the full pool.
% The proportion of false positives in the full negative pool is exactly
% $\mathrm{FAR}_{\text{base}} = FP_{\text{base}} / N_{\text{neg}}$.
% Applying this proportional sampling rule, the expected value
% $\mathbb{E}[\,\cdot\,]$ --- the long-run average count over many random
% outcomes --- of $a$ is:

% \[ \mathbb{E}[a] = N_{\text{skip\_neg}} \times \mathrm{FAR}_{\text{base}}
% = N_{\text{skip\_neg}} \times \frac{FP_{\text{base}}}{N_{\text{neg}}} \]

% We now compute the expected value of $F(\theta)$.
% Since $FP_{\text{base}}$ is a fixed constant determined by the baseline
% evaluation and does not vary randomly, it passes through the expectation
% unchanged:

% \[ \mathbb{E}[F(\theta)]
% = \frac{\mathbb{E}[a]}{FP_{\text{base}}}
% = \frac{N_{\text{skip\_neg}} \times \dfrac{FP_{\text{base}}}{N_{\text{neg}}}}{FP_{\text{base}}}
% = \frac{N_{\text{skip\_neg}}}{N_{\text{neg}}}
% = S_r(\theta) \]

% The $FP_{\text{base}}$ terms cancel exactly, and the result is the skip rate
% $S_r(\theta)$. Therefore, $F(\theta)$ carries no additional information beyond
% what $S_r(\theta)$ already captures: any difference between them in a specific
% experimental run is random variation, not a systematic signal attributable to
% the hyperparameter configuration $\theta$. Including both terms in the scoring
% objective would implicitly assign skip-rate behaviour a combined weight of
% $w_S + w_F$, while creating the appearance of a balanced objective with three
% independent criteria. We therefore exclude $F(\theta)$ from the scoring
% criterion and retain $\mathrm{FAR}(\theta)$ solely as an observed consequence
% metric reported in the results tables.